\let\labelindent\relax

%=====================================================
% Packages
%=====================================================
\usepackage{balance} % To balance the references
\usepackage[dvipsnames,table]{xcolor} % To have text colouring options
\usepackage{url} % To have links
\usepackage{graphicx} % Additional options for includegraphics
\usepackage{todonotes} % To have TODOs and missingfigures
\usepackage{booktabs} % needed for nice tables (toprule, bottomrule)
\usepackage{bbm} % indicator function
\usepackage{mathtools}
\usepackage{etoolbox} % robustify command

\usepackage{amsmath}
\usepackage{amsfonts}
\usepackage{pifont} % math symbols
\usepackage{amssymb} % math symbols

\usepackage{tikz} % create figures with tikz

\usepackage{hyperref} % to show links to the references, and also back to where the reference was cited

\usepackage{lipsum} % Dummy text (lorem ipsum)
\usepackage[mode=match]{siunitx} % SI units

% Notes in tables
\usepackage{threeparttable}
\usepackage{tablefootnote}

\usepackage{enumitem} % Customize itemize

%=====================================================
% Setup link colours (using hyperref)
%=====================================================

% Colors here (dvipsnames): https://www.overleaf.com/learn/latex/Using_colors_in_LaTeX

\hypersetup{
colorlinks=true
,linkcolor=NavyBlue
,citecolor=NavyBlue
% ,filecolor=MK_Two_Six
,urlcolor=NavyBlue
% ,menucolor=MK_Two_Five
% ,runcolor=MK_Two_Four
% ,linkbordercolor=MK_Two_One
,citebordercolor=NavyBlue
% ,filebordercolor=MK_Two_Three
% ,urlbordercolor=MK_Two_Six
% ,menubordercolor=MK_Two_Five
% ,runbordercolor=MK_Two_Four
}

%=====================================================
% Setup of SI Units
%=====================================================
\sisetup{per-mode = symbol,
         detect-weight = true,
         range-phrase = --,
         range-units = single,
         detect-all = true}

%=====================================================
% Commands
%=====================================================
\newcommand{\ra}[1]{\renewcommand{\arraystretch}{#1}}

% Inter-reference helpers (for figures, sections, etc)
\def\secref#1{Sec.~\ref{#1}}
\def\appref#1{Appendix ~\ref{#1}}
\def\figref#1{Fig.~\ref{#1}}
\def\tabref#1{Tab.~\ref{#1}}
\def\eqref#1{Eq.~(\ref{#1})}
\def\algref#1{Alg.~\ref{#1}}

% Random commands
% to write etal with the right format
\newcommand{\etal}{et al.}

% To write text vertically
\newcommand{\vt}[1]{\parbox[t]{2mm}{\multirow{3}{*}{\rotatebox[origin=c]{90}{#1}}}}

%=====================================================
% Extra symbols
%=====================================================
\newcommand{\cmark}{\ding{51}}%
\newcommand{\xmark}{\ding{55}}%

\newcommand{\magentasquare}{{\textcolor{magenta}{$\blacksquare$}}}
\newcommand{\bluesquare}{{\textcolor{blue}{$\blacksquare$}}}

\usepackage{tikz}
\newcommand*\notecircle[1]{\tikz[baseline=(char.base)]{
		\node[shape=circle,draw,inner sep=1pt, thick] (char) 
		{\footnotesize{#1}};}}
\newcommand{\boldcircled}[1]{\notecircle{\textbf{#1}}}

\newcommand*\notesquared[1]{\tikz[baseline=(char.base)]{
		\node[shape=rectangle,draw,inner sep=2pt, thick] (char) 
		{\footnotesize{#1}};}}
\newcommand{\boldsquared}[1]{\notesquared{\textbf{#1}}}

%=====================================================
% Table column types
% https://tex.stackexchange.com/a/12712
%=====================================================
\newcolumntype{L}[1]{>{\raggedright\let\newline\\\arraybackslash\hspace{0pt}}m{#1}}
\newcolumntype{C}[1]{>{\centering\let\newline\\\arraybackslash\hspace{0pt}}m{#1}}
\newcolumntype{R}[1]{>{\raggedleft\let\newline\\\arraybackslash\hspace{0pt}}m{#1}}
